\section{Discussion and policy implications}
The strands share a question, not a unit of observation. The survey distinguishes two reported uses. The interviews show how some women combine wallets, cash, trusted contacts and help from others. The expert findings describe why a product can work in one setting while another part of a transaction remains difficult. The benchmark documents what selected systems are designed to do. The strands complement each other without estimating a mechanism or validating one another's counts.

Three practical tests follow. First, can customers open and continue using an account under their own name when an accepted document changes, and get help when verification fails? Second, can a user see the total price and the amount due at destination before confirming a transfer? Third, if a transfer spans providers or an intermediary, can the user identify who handles the complaint, trace the payment and obtain a remedy? These tests separate usable access, transparent pricing and recourse from the number or speed of transactions. None presumes one universal barrier or treats cash use as failure.

The welfare implications concern things worth measuring: realized fees and exchange-rate spreads, travel or waiting time, delivery, losses and recovery, privacy, and control over spending. Some of these issues appear in the authored accounts; the survey's binary use outcomes cannot measure gains in them. Estimating welfare changes or ranking reforms would require outcomes and comparisons the present study does not provide. The contribution is to show why these questions should be asked at different stages of a payment.
